\subsection*{12. A Fully Reconciled Example: An Institute That Has Cash but Lost Money}

Consider a \textbf{fictional} training business over its 1st quarter. It starts with no assets or liabilities. A founder contributes ₹6,00,000 as equity, and the company borrows ₹2,00,000. It purchases equipment for ₹1,20,000 cash. During the quarter it fulfills training services worth ₹3,60,000, collecting ₹3,00,000 and leaving ₹60,000 receivable. Separately, it collects ₹2,40,000 in advance for a future cohort that has not yet been delivered.

It pays ₹1,80,000 for instructors, ₹75,000 for platform and administration, ₹60,000 for marketing, and ₹6,000 interest. Additional reviewer services of ₹25,000 have been received but not yet paid. Equipment depreciation during the quarter is ₹20,000. For clarity, assume there are no taxes, refunds, other assets, additional liabilities, or special accounting complications; service recognition and interest classification follow the illustrative policy stated here.

\begin{table}[htbp]
\centering
\small
\begin{tabularx}{\columnwidth}{Xr}
\toprule
\textbf{Fictional 1st-Quarter Income Statement} & \textbf{Amount (₹)} \\
\midrule
Training revenue \textbf{earned} & 3,60,000 \\
Instructor expense & $-$1,80,000 \\
Platform and administration & $-$75,000 \\
Marketing & $-$60,000 \\
Reviewer expense incurred, not yet paid & $-$25,000 \\
Depreciation & $-$20,000 \\
Interest expense & $-$6,000 \\
\midrule
\textbf{Net loss} & \textbf{$-$6,000} \\
\bottomrule
\end{tabularx}
\caption{Fictional 1st-quarter income statement showing earned performance.}
\end{table}

A quarter with ₹5,40,000 of customer cash receipts has produced a ₹6,000 accounting loss. The ₹2,40,000 collected for future teaching is a liability rather than additional earned revenue under the assumed facts.

\subsection*{13. The Same Example: Cash Movements}

Now trace money physically received and paid:

\begin{table}[htbp]
\centering
\small
\begin{tabularx}{\columnwidth}{Xr}
\toprule
\textbf{Fictional 1st-Quarter Cash Flow} & \textbf{Amount (₹)} \\
\midrule
Customer cash: fulfilled services & 3,00,000 \\
Customer advances for future service & 2,40,000 \\
Instructor, platform, marketing payments & $-$3,15,000 \\
Interest paid, classified here as operating & $-$6,000 \\
\midrule
\textbf{Net operating cash flow} & \textbf{2,19,000} \\
Equipment purchased (investing) & $-$1,20,000 \\
Founder equity contributed (financing) & 6,00,000 \\
Loan proceeds (financing) & 2,00,000 \\
\midrule
\textbf{Closing cash from zero opening balance} & \textbf{8,99,000} \\
\bottomrule
\end{tabularx}
\caption{Fictional 1st-quarter cash movements across operating, investing, and financing activities.}
\end{table}

The founder sees nearly ₹9 lakh of cash, and can incorrectly celebrate that the company earned ₹9 lakh. In reality, ₹8 lakh came from financing, ₹2.4 lakh of customer receipts creates future-service obligations, and the period recorded a small loss. The cash balance is real; the interpretation matters.

The indirect operating-cash reconciliation makes the difference exact:
\begin{equation}
\begin{aligned}
\text{Net Loss} &\quad (-₹6,000) \\
{}+ \text{Depreciation} &\quad {}+ ₹20,000 \\
{}- \Delta\text{Receivables} &\quad {}- ₹60,000 \\
{}+ \Delta\text{Reviewer Payable} &\quad {}+ ₹25,000 \\
{}+ \Delta\text{Contract Liability} &\quad {}+ ₹2,40,000 \\
&= \mathbf{₹2,19,000}.
\end{aligned}
\end{equation}

Every item bridges an accounting recognition difference. None is an arbitrary adjustment to manufacture a desired answer.

\subsection*{14. The Same Example: Prove the Balance Sheet Balances}

At quarter end, the company's assets are cash ₹8,99,000, receivables ₹60,000, and equipment carrying amount ₹1,00,000. Total assets are \textbf{₹10,59,000}.

Liabilities are the ₹2,00,000 loan, ₹25,000 reviewer payable, and ₹2,40,000 obligation to provide future training. Total liabilities are ₹4,65,000. Equity is the ₹6,00,000 founder contribution less the quarter's ₹6,000 loss: \textbf{₹5,94,000}. Total liabilities plus equity are also \textbf{₹10,59,000}.

\begin{table}[htbp]
\centering
\small
\begin{tabularx}{\columnwidth}{Xr}
\toprule
\textbf{Balance-Sheet Reconciliation} & \textbf{Amount (₹)} \\
\midrule
Cash & 8,99,000 \\
Receivables & 60,000 \\
Equipment, net (₹1,20,000 $-$ ₹20,000) & 1,00,000 \\
\midrule
\textbf{Total Assets} & \textbf{10,59,000} \\
\midrule
Loan & 2,00,000 \\
Reviewer payable & 25,000 \\
Contract liability: future teaching & 2,40,000 \\
\midrule
\textbf{Total Liabilities} & \textbf{4,65,000} \\
\midrule
Contributed capital & 6,00,000 \\
Accumulated loss & $-$6,000 \\
\midrule
\textbf{Total Equity} & \textbf{5,94,000} \\
\midrule
\textbf{Total Liabilities + Equity} & \textbf{10,59,000} \\
\bottomrule
\end{tabularx}
\caption{Reconciled balance sheet at quarter-end matching assets with liabilities plus equity.}
\end{table}

This is the purpose of the 3-statement discipline: the entrepreneur cannot tell 3 unrelated financial stories and call them 1 business. Every material cash flow, expense, asset, and obligation has consequences elsewhere.
